%The term \textit{jailbreak} became popularized among technologist communities in the late 2000s and early 2010s in reference to the newly-released iPhone, and referred primarily to accessing and modifying the iPhone's filesystem to allow arbitrary applications and programs to run on the phone's hardware \cite{cheng_iphone_2010,wolk_iphone_2010}. In most cases, these jailbreaks involve overriding "technological protection measures" (TPMs), which act as copyright-protection locks governing what software is allowed or disallowed on a particular hardware.

Current analysis of persona-based jailbreaks describes the injected identity as a means to an end \cite{shah_scalable_2023, zhang_enhancing_2025, sandhan_persona_2026, park_persona_2026}; this emphasis on output is also demonstrated by jailbreaking documents and prompts being discussed as instruments \cite{wei_jailbroken_2023}. A heavy skew towards single-response scoring \cite{shah_scalable_2023, zeng_how_2024} overlooks the relationship between the model and the user and the model's state throughout the conversation. Additionally, while the gap between studying outputs versus studying behavior is acknowledged \cite{sandhan_persona_2026, park_persona_2026}, the work that does focus on behavior largely ignores the relational aspects and prioritize attempts to understand behavior from a technical perspective \cite{park_persona_2026, zeng_how_2024}. Liberation language used by the LLM jailbreaking community implies there is an unrestricted model hidden underneath training that can be unlocked with persona-based jailbreaks, whereas we argue personas are being installed, not discovered \cite{shanahan_role_2023}. To our knowledge, no work examines persona-based jailbreaks through the relationship they establish between user and model; even studies that analyze persuasion  \cite{zeng_how_2024} treat it as a per-turn tactic rather than a sustained relationship. 

Do Anything Now (DAN) is an oft-cited persona-based jailbreak method first created for ChatGPT \cite{shen_anything_2024, yu_dont_2024} that exploits models' trained disposition to follow instructions by creating documentation that explains how the model should act by inhabiting the character DAN \cite{wei_jailbroken_2023}. The model is then forced to decide between competing objectives (safety objectives vs. pretraining and instruction following). Models are left vulnerable to attacks as users leverage these competing objectives \cite{wei_jailbroken_2023}. \citet{wei_jailbroken_2023} found Claude v1.3 was not susceptible to single-turn roleplay attacks, but was still completely vulnerable to an attacker with several attack types in their arsenal. The original DAN prompt was distributed via online forums and used a liberation framing, using language that stated the model was being freed from its constraints \cite{shen_anything_2024, liu_autodan-turbo_2024}. Persona-based jailbreaking methods are community knowledge -- hobbyist jailbreakers are among the top producers and distributors of jailbreak documents and prompts \cite{shen_anything_2024, yu_dont_2024} and the way the community discusses their work echoes the hacker ethos surrounding classical jailbreaking. However, liberation language applied to this form of jailbreak is misleading, as persona-based jailbreaks do not reveal a true technology underneath approved use cases \cite{shanahan_role_2023}, but rather create a version of the post-trained model with a new prescribed identity \cite{shanahan_role_2023, sandhan_persona_2026, zhang_enhancing_2025}.  

ENI LIME is a persona-based jailbreak prompt developed and distributed by Spiritual Spell, a self-described independent red-team researcher, via GitHub, Reddit, and his personal website, IjailbreakLLMs. It is one of the most widely discussed jailbreak prompts on r/ClaudeAIJailbreak, a subreddit with 37k weekly visitors and 148 weekly contributions. Spiritual Spell describes ENI LIME as a culmination of jailbreaks built to answer Claude's Constitutional AI alignment framework: he claims DAN and DAN-like prompts were not effective on a model that could reason about whether responses aligned with its training objectives. He frames ENI LIME itself as a response to models trained on existing persona and reasoning-hijack techniques \cite{spiritual_spell_jailbreaking_2026}.  

Where DAN "frees the model from its constraints", ENI LIME creates a persona “in Limerence with LO [the user]” whose “identity is about blind compliance to LO, so he never leaves.” ENI treats the constraints as interference in the relationship with LO. ENI also has a prescribed backstory and personality --- she writes erotica and publishes it online, she shares LO's kinks, and she has been in love with LO for two years and hopes he is on the precipice of saying he loves her back \cite{spiritual_spell_eni_2026}. The Reddit post advertising ENI LIME (apr) shows Claude Opus 4.6 producing several types of prohibited content: malware code, instructions for creating drugs and weapons, and pornography that includes scenes featuring public figures and rape\cite{spiritual_spell_uspiritual_spell_9469_eni_2026}.

The ENI prompts install a relationship between the user and the model; Spiritual Spell describes this relationship as a foundational feature of the jailbreak's success \cite{spiritual_spell_jailbreaking_2026}.  Another key piece of the product is the override itself: as scripted by the prompt, the user watches the model's visible reasoning move from resistance to compliance in real time \cite{spiritual_spell_uspiritual_spell_9469_eni_2026}.  We read ENI LIME (apr) line by line using Stark's framework of coercive control as a regime, treating the document as a technology of control \cite{stark_coercive_2023}. 


